\paragraph{两个\trFull\trEmbedding{}。}\label{sub:two-full-embeddings}
此时，引进一个技术性条件是方便的。称 $\pi\colon E\to S$ 具有\textbf{\trGeneric\trParallelTransport}（generic parallel transport），若在\trVanishingPath{}的\trDistinguishedBasis{}的每一个\trHomotopy{}类之内，都存在一个其相伴的\trVanishingCycle{}处于\trGeneralPosition{}的基。给定任何一个\trLefschetzFibration，都可以找到 $(\theta_E,\phi_E,I_E)$ 的一个小\trPerturbation，它在 $\partial E$ 附近并且在 $M=E_*$ 附近都是常值的，使得\trPerturbation{}后的结构具有\trGeneric\trParallelTransport（这容易由\cref{thm:transport-isotopy-realization}以及\trDistinguishedBasis{}只有可数多个\trHomotopy{}类这一事实推出）。这样的\trPerturbation{}不改变 $M$ 或 $E$ 的\trFukayaCategory（在第二种情形，这是在\trQuasiisomorphic{}的意义下，并且要用到\trMoser{}引理）。记住了这一点，从现在起我们将假定，我们给定的\trLefschetzFibration{}已经具有\trGeneric\trParallelTransport。

从现在起我们要求 $\mathrm{char}(\KK)\neq 2$；这使我们可以应用 §6 中的结果，而更重要的是，它使得\cref{thm:type-b-equivariant-hf}中的 $\mathrm{Id}$ 不同于 $-\mathrm{Id}$。取 $\widetilde{\scr{A}}\subset\scr{F}(\widetilde{E})^{\mathrm{naive}}$ 为那样的\trFull{} \trAinfty-子\trCategory：其\trObject{}是诸预\trEquivariant{}膜 $\widetilde{L}^{\#}_{\pm}$ 以及诸\trEquivariant{}膜 $L^{\#}$、$\widetilde{\Delta}^{\#}_\gamma$。尽管这在很大程度上重复了前面的讨论，我们仍然想明确说明，在构造 $\widetilde{\scr{A}}$ 时用到的是哪些种类的\trFloerDatum{}与\trPerturbationDatum。对任何一对 (U) 型\trEquivariant{}或预\trEquivariant{}膜，\trFloerDatum{}的选取应使 \eqref{eq:tilde-cf-decomposition} 成立。这特别地蕴含：每个这样的膜都是 $H^0(\widetilde{\scr{A}})$ 的一个自由\trObject，其意义按\cref{sub:z2-basic-notions}理解。类似地，对任何由 (U) 型膜组成的 $(d{+}1)$ 元组（$d\ge 2$），诸\trPerturbationDatum{}也取同一类型的，以使该\trIsomorphism{}与诸 \trAinfty-\trStructureMap{}\trCompatible。其次，对任何一对 (B) 型\trEquivariant{}膜，\trFloerDatum{}应当是满足\cref{thm:floer-datum-assumption}的一个不变数据。然而，在诸膜来自使 \eqref{eq:two-vanishing-paths-position} 成立的路径 $\gamma_0$、$\gamma_1$ 的特殊情形，我们要采用导致\trIsomorphism{}\eqref{eq:cf-delta-vanishing-iso}的那个更严格的选取。每当由\trVanishingPath{}组成的一个集合 $(\gamma_0,\ldots,\gamma_d)$ 满足 \eqref{eq:d-vanishing-paths-position} 时，对诸\trPerturbationDatum{}的选取也将适用同样的规则。这与对全局选取的诸\trConsistency{}条件是\trCompatible{}的，因为这类 $(d{+}1)$ 元组在取有序子集之下是封闭的。
\begin{definition}\label{thm:fukaya-category-lefschetz-fibration}
\trLefschetzFibration{}的\trFukayaCategory，记作 $\scr{F}(\pi)$，是不变部分 $\widetilde{\scr{A}}^{\mathrm{inv}}$。
\end{definition}
\begin{proposition}\label{thm:full-embedding-functor}
存在一个\trCohomology{}\trFullyFaithful{}的 \trAFunctor
\begin{equation*}
\scr{F}(E)\longrightarrow\scr{F}(\pi),
\end{equation*}
它把每个位于 $E^{\mathrm{shrink}}$ 之内的\trExactLagrangianBrane{}$L^{\#}$ 送到其\trPreimage{}$\widetilde{L}^{\#}$。
\end{proposition}
\begin{proof}
按假定，$E$ 上的每个\trExactLagrangianBrane{}都 \trIsotopy{}于一个位于 $E^{\mathrm{shrink}}$ 之内的膜。因此，由这类膜组成的\trFull{} \trAinfty-子\trCategory{}到 $\scr{F}(E)$ 中的包含是一个\trQuasiequivalence，而它可以由\cref{thm:quasi-equivalence-homotopy-inverse}求逆（事实上，只要看一下紧接该定理之前的论证，就能找到一个在子\trCategory{}自身上为恒同的逆\trQuasiequivalence）。另一方面，同一个子\trCategory{}又可以借 \eqref{eq:tilde-cf-decomposition} 与 $\widetilde{\scr{A}}^{\mathrm{inv}}$ 中由诸\trObject{}$\widetilde{L}^{\#}$ 组成的那个\trFull{} \trAinfty-子\trCategory{}相认同。
\end{proof}
现在设 $\boldsymbol{\gamma}=(\gamma_1,\ldots,\gamma_m)$ 是\trVanishingPath{}的一个\trDistinguishedBasis。对每个 $k=1,\ldots,m$，选好一个\trGrading，使与之相伴的 (B) 型\trLagrange\trSubmanifold{}$\widetilde{\Delta}_k$ 成为一个\trEquivariant{}膜，并给相应的\trVanishingCycle{}$V_k$ 配备诱导的\trBraneStructure。把如此所得的 $\scr{F}(\pi)$ 与 $\scr{F}(M)$ 中\trObject{}的两个集合分别记作 $\widetilde{\boldsymbol{\Delta}}^{\#}$ 与 $\boldsymbol{V}^{\#}$。它们各自带有\trDirected{} \trAinfty-子\trCategory{}$\scr{F}^{\rightarrow}(\widetilde{\boldsymbol{\Delta}}^{\#})$ 与 $\scr{F}^{\rightarrow}(\boldsymbol{V}^{\#})$。
\begin{proposition}\label{thm:fibration-exceptional-collection}
$\widetilde{\boldsymbol{\Delta}}^{\#}$ 是 $\scr{F}(\pi)$ 中的一个\trExceptionalCollection。此外，$\scr{F}^{\rightarrow}(\widetilde{\boldsymbol{\Delta}}^{\#})$ 与 $\scr{F}^{\rightarrow}(\boldsymbol{V}^{\#})$ 是\trQuasiisomorphic{}的，所经的\trAFunctor{}是\trStrictlyUnital{}的，并且保持诸\trObject{}的排序。作为推论，存在一个\trCohomology{}\trFullyFaithful{}的 \trAFunctor
\begin{equation*}
\scr{F}^{\rightarrow}(\boldsymbol{V}^{\#})\longrightarrow\scr{F}(\pi),
\end{equation*}
它把每个 $V^{\#}_k$ 送到 $\widetilde{\Delta}^{\#}_k$。
\end{proposition}
\begin{proof}
先设诸\trVanishingCycle{}处于\trGeneralPosition{}。由\cref{thm:type-b-equivariant-hf}(i) 我们知道
\begin{equation*}
\Hom_{H\scr{F}(\pi)}(\widetilde{\Delta}^{\#}_k,\widetilde{\Delta}^{\#}_k)
=HF^*(\widetilde{\Delta}^{\#}_k,\widetilde{\Delta}^{\#}_k)^{\ZZ/2}
\cong H^*(S^{n+1};\KK)^{\ZZ/2}=\KK,
\end{equation*}
因此每个 $\widetilde{\Delta}^{\#}_k$ 都是一个例外\trObject。此外，一如\cref{thm:dual-action-minus-id}中已经提到的，同一条引理的第 (ii) 部分与 \eqref{eq:cf-delta-vanishing-iso} 一起蕴含：当 $j>k$ 时，$\ZZ/2$-作用在 $HF^*(\widetilde{\Delta}^{\#}_j,\widetilde{\Delta}^{\#}_k)$ 上是 $-\mathrm{Id}$，从而 $\scr{F}(\pi)$ 中的\trMorphism{}群是\trAcyclic{}的，这就证明了我们确实有一个\trExceptionalCollection。\cref{thm:directed-subcategory-inclusion}随后便提供一个从 $\scr{F}^{\rightarrow}(\widetilde{\boldsymbol{\Delta}}^{\#})$ 到 $\scr{F}(\pi)$ 的\trCohomology{}\trFullyFaithful{}\trFunctor。另一方面，由 \eqref{eq:cf-delta-vanishing-iso} 可知，那个\trDirected{}子\trCategory{}与 $\scr{F}^{\rightarrow}(\boldsymbol{V}^{\#})$ 相同。
对于一般情形，我们利用\trLefschetzFibration{}具有\trGeneric\trParallelTransport{}这一假定，它蕴含：对任何 $\boldsymbol{\gamma}$，都有一个与之\trHomotopic{}的 $\boldsymbol{\gamma}'$ 使前面的论证适用。由于每个\trVanishingCycle{}$V'^{\,n}_k$ 都 \trIsotopy{}于 $V^n_k$，它们在 $H^0\scr{F}(M)$ 中成为\trIsomorphic{}的\trObject。\cref{thm:directed-collections-quasi-isomorphic}表明，由此所得的两个\trDirected{} \trAinfty-子\trCategory{}是\trQuasiisomorphic{}的。类似地，每个 $\widetilde{\Delta}'_k$ 都经由 $\ZZ/2$-不变的\trLagrange\trSubmanifold{} \trIsotopy{}于 $\widetilde{\Delta}_k$，从而 $H^0\scr{F}(\pi)$ 中与之相伴的\trObject{}是\trIsomorphic{}的，这就给出了介于与之相伴的（而且也是\trFull{}的）\trDirected{} \trAinfty-子\trCategory{}之间的一个\trQuasiisomorphism。
\end{proof}
